\chapter{Objetivos, alcance y requisitos actualizados}

\section{Objetivo general}
El objetivo general no ha cambiado respecto al GEP: diseñar y desplegar un asistente personal local, autoalojado y proactivo que combine inteligencia artificial, memoria documental, domótica y operación segura bajo control del usuario. La formulación actual es más precisa porque el proyecto ya ha dejado de ser hipotético. Odín no pretende construir un modelo fundacional propio, sino una arquitectura capaz de hacer útil la IA local sobre datos y servicios personales reales.

\section{Objetivos específicos revisados}
\begin{longtable}{p{0.08\textwidth}p{0.34\textwidth}p{0.16\textwidth}p{0.34\textwidth}}
  \caption{Objetivos específicos revisados}
  \label{tab:objetivos-revisados}\\
  \toprule
  ID & Objetivo & Estado & Evidencia actual \\
  \midrule
  \endfirsthead
  \toprule
  ID & Objetivo & Estado & Evidencia actual \\
  \midrule
  \endhead
  \bottomrule
  \endfoot
  O1 & Desplegar infraestructura propia reproducible. & Avanzado & Servidor, Docker Compose, carpetas por dominio y repositorio saneado. \\
  O2 & Ejecutar inferencia local sin depender de APIs externas. & Avanzado & Ollama operativo como servicio del sistema. \\
  O3 & Construir memoria personal consultable. & Avanzado & Nextcloud, Qdrant, ingesta y búsqueda documental. \\
  O4 & Integrar conversación con acciones reales. & Avanzado & Once herramientas de Open WebUI y webhooks de n8n. \\
  O5 & Controlar el entorno físico. & Avanzado & Home Assistant, HA Voice, presencia, luces, música, temporizadores y cámaras. \\
  O6 & Mantener acceso remoto seguro. & Avanzado & Tailscale adoptado como VPN principal. \\
  O7 & Incorporar monitorización y proactividad. & Avanzado & Netdata, Telegram, exportador HA y autoreparabilidad conservadora. \\
  O8 & Definir estrategia de backup y continuidad. & Parcial & Política con NVMe externo definida; falta prueba final de restauración. \\
  O9 & Documentar decisiones, límites y resultados. & En curso & Memoria extensa, bibliografía real y anexos técnicos saneados. \\
  \bottomrule
\end{longtable}

\section{Correcciones respecto al alcance inicial}
El GEP inicial era deliberadamente ambicioso y hablaba de un ``asistente total''. El seguimiento permite separar mejor tres categorías:
\begin{itemize}
  \item \textbf{núcleo obligatorio}: memoria, conversación, domótica, fotos, acceso remoto y operación segura;
  \item \textbf{capacidades ya alcanzadas}: notas, búsqueda documental, Immich, Frigate, recordatorios, lista de compra, salud del sistema y alertas;
  \item \textbf{líneas pospuestas o fuera de alcance}: clasificación avanzada de correos, avisos de contraseñas, dashboard único terminado y agente de escritorio.
\end{itemize}

Esta corrección mejora el rigor del proyecto. No todo lo imaginable debe entrar en el TFG para que el sistema sea valioso. De hecho, decidir que correos y contraseñas quedan fuera de alcance evita duplicar servicios especializados que ya resuelven bien esos problemas y permite concentrar el esfuerzo en lo diferencial de Odín.

\section{Requisitos funcionales}
\begin{longtable}{p{0.11\textwidth}p{0.39\textwidth}p{0.38\textwidth}}
  \caption{Requisitos funcionales del seguimiento}
  \label{tab:requisitos-funcionales-seguimiento}\\
  \toprule
  Código & Requisito & Implementación o evidencia \\
  \midrule
  \endfirsthead
  \toprule
  Código & Requisito & Implementación o evidencia \\
  \midrule
  \endhead
  \bottomrule
  \endfoot
  RF1 & Guardar notas, ideas o documentos desde lenguaje natural. & Tool de guardado, n8n y Nextcloud. \\
  RF2 & Buscar información privada en documentos. & Nextcloud, Qdrant y capa de recuperación documental. \\
  RF3 & Buscar recuerdos visuales. & Immich conectado al chat. \\
  RF4 & Controlar dispositivos domóticos. & Home Assistant y herramientas conversacionales. \\
  RF5 & Crear recordatorios y consultar listas. & Calendario y lista de la compra de Home Assistant. \\
  RF6 & Consultar cámaras y eventos recientes. & Frigate integrado en conversación. \\
  RF7 & Consultar salud del servidor. & Tool de salud, Netdata y exportador JSON. \\
  RF8 & Avisar de incidencias relevantes. & Telegram diario y alertas reactivas. \\
  RF9 & Capturar contenido web o vídeo. & Transcripciones de YouTube y lectura web. \\
  RF10 & Acceder de forma remota segura. & Tailscale sobre red privada. \\
  \bottomrule
\end{longtable}

\section{Requisitos no funcionales}
\begin{table}[H]
  \centering
  \caption{Requisitos no funcionales}
  \label{tab:requisitos-no-funcionales}
  \begin{tabularx}{\textwidth}{lX}
    \toprule
    Requisito & Interpretación en Odín \\
    \midrule
    Privacidad & Priorizar ejecución local, no publicar secretos y minimizar exposición de servicios. \\
    Mantenibilidad & Usar Compose, documentación versionada y separación progresiva por dominios. \\
    Disponibilidad & Mantener Home Assistant desacoplado del servidor principal y monitorizar servicios. \\
    Extensibilidad & Añadir nuevas tools sin rediseñar toda la arquitectura. \\
    Reproducibilidad & Conservar configuraciones saneadas y decisiones técnicas trazables. \\
    Usabilidad & Preferir flujos conversacionales y respuestas realmente útiles frente a demos aisladas. \\
    \bottomrule
  \end{tabularx}
\end{table}

\section{Dispositivos opcionales y principio de integración}
Una de las ideas más importantes del proyecto es que Odín no exige un catálogo cerrado de dispositivos. La infraestructura mínima puede funcionar solo con el servidor y la capa software. A partir de ahí, Home Assistant permite incorporar dispositivos ya existentes en casa, incluso si no fueron adquiridos específicamente para el TFG.

\begin{table}[H]
  \centering
  \caption{Dispositivos físicos opcionales integrables}
  \label{tab:dispositivos-opcionales}
  \begin{tabularx}{\textwidth}{lXX}
    \toprule
    Dispositivo & Papel posible & Tratamiento económico \\
    \midrule
    Sensor de presencia & Detectar llegada o actividad y disparar automatizaciones. & Opcional; coste incremental cero si ya estaba disponible. \\
    Enchufe inteligente & Encender, apagar o medir consumo de equipos concretos. & Opcional; útil para futuras métricas energéticas. \\
    Luces inteligentes & Control conversacional de la habitación. & Opcional; integrables como entidades HA. \\
    Altavoz o satélite de voz & Interacción oral distribuida. & Opcional; HA Voice ya cubre el caso principal. \\
    Cámara compatible & Eventos y capturas para Frigate. & Opcional; amplía seguridad y contexto visual. \\
    \bottomrule
  \end{tabularx}
\end{table}

La gracia arquitectónica es que el modelo no necesita conocer de fábrica cada marca ni cada dispositivo. Basta con exponer las entidades a Home Assistant y describir de forma clara qué son y cómo se interactúa con ellas. La IA puede razonar sobre esas capacidades y convertir peticiones naturales en acciones concretas. Esto hace que Odín sea ampliable sin quedar atado a un ecosistema propietario ni a una lista cerrada de hardware.
